\section{Inter-node performance}
\label{section:tailoring:inter-node}

% not a huge amount of content here, just a note perhaps on how we can scale up/out are conclusions on the intra-node level thread/rank placement/affinity, e.g., how likwid interprets your directions on stating: "give me 8 ranks, 1 per node" -> it will then know to run on 8 nodes

% there is perhaps more we could say for environment variables or MPI runtime, but I always get confused

% finally, are there anythings we need consider with the network? I presume this is all fixed, but being able to pick the nodes we run on could help?

At the inter-node level we are now more concerned about our use of the network.
HPC systems can have variability of interconnect speeds between nodes, and some
compute centres may make this data available hence users can select specific
nodes through a scheduler like SLURM to avoid poorly performing connections.
However, this is a highly system specific approach and so we do not consider
this issue here.

Firstly, we need to remember that we can scale up our LIKWID to a multi-node configuration...

Secondly, \texttt{#SBATCH --contiguous} and \texttt{#SBATCH --switches=1}...

Thirdly, anything from \texttt{mpirun}...?
